\documentclass[a4paper, serif, 10pt]{moderncv}

%% moderncv
% style and color
\moderncvstyle{banking}
\moderncvcolor{black}

% renew moderncv command to adapt for CJK colon
\renewcommand*{\cvitem}[3][.25em]{%
  \ifstrempty{#2}{}{\hintstyle{#2}：}{#3}%
  \par\addvspace{#1}}

\renewcommand*{\cvitemwithcomment}[4][.25em]{%
  \savebox{\cvitemwithcommentmainbox}{\ifstrempty{#2}{}{\hintstyle{#2}：}#3}%
  \setlength{\cvitemwithcommentmainlength}{\widthof{\usebox{\cvitemwithcommentmainbox}}}%
  \setlength{\cvitemwithcommentcommentlength}{\maincolumnwidth-\separatorcolumnwidth-\cvitemwithcommentmainlength}%
  \begin{minipage}[t]{\cvitemwithcommentmainlength}\usebox{\cvitemwithcommentmainbox}\end{minipage}%
  \hfill% fill of \separatorcolumnwidth
  \begin{minipage}[t]{\cvitemwithcommentcommentlength}\raggedleft\small\itshape#4\end{minipage}%
  \par\addvspace{#1}}

%% page layout/margins
\usepackage[top=2.5cm, bottom=2.5cm, left=1.5cm, right=1.5cm]{geometry}

\nopagenumbers{}

%% Babel config for English language
\usepackage[english]{babel}

%% fontspec
\usepackage{fontspec}

\IfFontExistsTF{Linux Libertine}{
  \setmainfont[Ligatures={TeX, Common}, Numbers=Lining, ItalicFont=Linux Libertine]{Linux Libertine}
}{}
\IfFontExistsTF{Linux Libertine O}{
  \setmainfont[Ligatures={TeX, Common}, Numbers=Lining, ItalicFont=Linux Libertine O]{Linux Libertine O}
}{}

%% CTeX
% CJK support, used to show CJK characters in the resume
%
% - fontset=none: disable builtin fontset but instead set the CJK font manually
% - heading=false: disable ctex heading
% - punct=kaiming: use kaiming punctuations styles for CJK
% - scheme=plain: use plain scheme, do not override \normalsize font size
% - space=auto: space settings for CJK characters
%
% ref:
% - http://ctan.mirrorcatalogs.com/language/chinese/ctex/ctex.pdf
\usepackage[UTF8, heading=false, punct=kaiming, scheme=plain, space=auto]{ctex}

\IfFontExistsTF{Noto Serif CJK SC}{
  \setCJKmainfont{Noto Serif CJK SC}
}{}
\IfFontExistsTF{Noto Sans CJK SC}{
  \setCJKsansfont{Noto Sans CJK SC}
}{}

%% line spacing
\usepackage{setspace}
\setstretch{1.125}

%% URL styling - use normal text instead of monospace
\urlstyle{same}

% Auto-underline all links
% Defer until \begin{document} so hyperref is already loaded
\AtBeginDocument{%
  \let\oldhref\href
  \renewcommand{\href}[2]{\oldhref{#1}{\underline{#2}}}%
}

%% Patch \httplink and \httpslink to handle full URLs
% The original moderncv macros always prepend http:// or https:// to URLs.
% This causes issues when the URL already has a protocol (e.g., https://example.com)
% resulting in malformed URLs like https://https://example.com.
% This patch checks if the URL already contains "://" and uses it directly if so.
% Note: We use \str_set:Nx (x-expansion) to expand macros like \@homepage first.
\ExplSyntaxOn
\str_new:N \g_yamlresume_url_str

\RenewDocumentCommand{\httpslink}{O{}m}{
  \str_set:Nx \g_yamlresume_url_str {#2}
  \str_if_in:NnTF \g_yamlresume_url_str {://}
    { \url{#2} }
    { \url{https://#2} }
}

\RenewDocumentCommand{\httplink}{O{}m}{
  \str_set:Nx \g_yamlresume_url_str {#2}
  \str_if_in:NnTF \g_yamlresume_url_str {://}
    { \url{#2} }
    { \url{http://#2} }
}
\ExplSyntaxOff

\name{林雅婷}{}
\title{資深行銷經理 | 品牌策略與數位成長}
\phone[mobile]{+886 912 345 678}
\email{yating.lin@example.com}
\homepage{https://www.yatinglin.com}

\address{中國臺灣，臺北市，臺北市，臺北市信義區松高路 88 號 12 樓，110}{}{}

\extrainfo{{\small \faLinkedin}\ \href{https://www.linkedin.com/in/yating-lin-marketing}{@yating-lin-marketing} {} {} {} • {} {} {} 
{\small \faMedium}\ \href{https://medium.com/@yating.marketing}{@yating.marketing} {} {} {} • {} {} {} 
{\small \faInstagram}\ \href{https://www.instagram.com/yating.marketing}{@yating.marketing}}

\begin{document}

\maketitle

\section{簡介}

\cvline{}{\begin{itemize}
\item 超過 8 年行銷經驗，專注品牌策略、數位成長與全通路行銷，擅長以數據驅動決策。
\item 曾帶領 6 人團隊管理新台幣 8,000 萬元年度預算，帶動電商營收年增 35\%。
\item 熟悉 Google Analytics 4、Meta Ads、CRM 與行銷自動化，具備跨部門協作與專案管理能力。
\item 擁有國際交換與跨文化溝通經驗，能以中、英、日語與不同市場團隊合作。
\end{itemize}}
\section{教育背景}

\cventry{2016年9月–2018年6月}
        {碩士，企業管理研究所，成績：4.1/4.3}
        {國立臺灣大學}
        {\url{https://www.ntu.edu.tw/}}
        {}
        {\begin{itemize}
\item 專攻品牌管理與數位行銷，論文探討社群媒體對消費者購買意願的影響
\item 擔任行銷研究助理，協助企業進行市場調查與數據分析
\item 獲選為商學院國際交換學生，赴荷蘭鹿特丹管理學院進修一學期
\end{itemize}
\textbf{課程}：行銷管理、消費者行為、品牌策略、數位行銷分析、市場研究、策略管理}

\cventry{2012年9月–2016年6月}
        {學士，廣告學系，成績：3.9/4.0}
        {國立政治大學}
        {\url{https://www.nccu.edu.tw/}}
        {}
        {\begin{itemize}
\item 奠定廣告、公關與整合行銷傳播的理論基礎
\item 畢業專題以「行動支付品牌年輕化策略」獲得系上最佳專題獎
\item 主辦校園行銷營，吸引超過 200 位學生參與
\end{itemize}
\textbf{課程}：廣告理論、公關策略、媒體規劃、行銷學、統計學、創意策略}
\section{工作經歷}

\cventry{2021年3月至今}
        {資深行銷經理}
        {富邦媒體科技股份有限公司}
        {\url{https://www.momoshop.com.tw/}}
        {}
        {\begin{itemize}
\item 負責台灣與東南亞市場的品牌策略、數位廣告與會員成長，管理年度行銷預算新台幣 8,000 萬元
\item 帶領 6 人行銷團隊，規劃並執行全通路 campaign，帶動電商營收年增 35\%
\item 建立數據驅動的行銷決策流程，整合 Google Analytics、Meta Ads 與 CRM，提升 ROAS 至 4.8
\item 主導品牌重塑專案，重新定位品牌訊息與視覺識別，使品牌好感度提升 22\%
\item 與產品、業務及數據團隊跨部門協作，推動個人化推薦與自動化行銷
\end{itemize}
\textbf{關鍵字}：品牌策略、數位廣告、會員成長、數據分析、團隊領導、全通路行銷}

\cventry{2018年7月–2021年2月}
        {行銷經理}
        {奧美整合行銷傳播集團}
        {\url{https://www.ogilvy.com/tw}}
        {}
        {\begin{itemize}
\item 服務科技、金融與快速消費品客戶，負責整合行銷傳播策略與專案管理
\item 規劃社群媒體內容與 KOL 合作，單一 campaign 觸及超過 500 萬人次
\item 管理客戶關係與提案簡報，成功爭取 3 個新客戶，年度合約金額成長 40\%
\item 分析市場趨勢與競品動態，提供客戶即時的策略建議
\end{itemize}
\textbf{關鍵字}：整合行銷、社群媒體、KOL 合作、客戶管理、提案簡報}

\cventry{2016年7月–2018年6月}
        {行銷專員}
        {台灣樂天市場股份有限公司}
        {\url{https://www.rakuten.com.tw/}}
        {}
        {\begin{itemize}
\item 執行電子報、EDM 與站內促銷活動，提升活躍會員數 18\%
\item 協助規劃年度檔期活動，包含雙 11、週年慶與品牌日
\item 運用 Google Analytics 追蹤流量與轉換，優化 landing page 成效
\end{itemize}
\textbf{關鍵字}：電商行銷、EDM、促銷活動、Google Analytics}
\section{語言}

\cvline{普通話}{母語或雙語水平 \hfill \textbf{關鍵字}：母語}
\cvline{英語}{完全專業水平 \hfill \textbf{關鍵字}：TOEIC 950、國際交換經驗}
\cvline{日語}{有限工作水平 \hfill \textbf{關鍵字}：JLPT N3}
\section{專業技能}

\cvline{數位行銷}{專家 \hfill \textbf{關鍵字}：SEO、SEM、Google Ads、Meta Ads、Google Analytics 4、LINE 官方帳號}
\cvline{品牌與內容策略}{高級 \hfill \textbf{關鍵字}：品牌定位、內容行銷、社群經營、公關溝通、文案撰寫}
\cvline{數據分析與 CRM}{高級 \hfill \textbf{關鍵字}：Looker Studio、SQL、CRM、A/B Testing、行銷自動化}
\cvline{專案與團隊管理}{專家 \hfill \textbf{關鍵字}：跨部門協作、預算管理、敏捷開發、領導激勵、供應商管理}
\section{榮譽}

\cventry{2023年12月}
        {台灣數位行銷協會}
        {年度最佳行銷專案獎}
        {}
        {}
        {以「全通路會員成長計畫」榮獲年度最佳行銷專案獎，該專案在 12 個月內提升會員回購率 28\%。}

\cventry{2018年6月}
        {國立臺灣大學管理學院}
        {傑出學生論文獎}
        {}
        {}
        {碩士論文《社群媒體內容對消費者購買意願之影響》獲選為傑出學生論文。}
\section{證書}

\cventry{2023年3月}
        {Google}
        {Google Analytics 4 Certification}
        {\url{https://skillshop.exceedlms.com/student/catalog}}
        {}
        {}

\cventry{2022年8月}
        {Meta}
        {Meta Certified Digital Marketing Associate}
        {\url{https://www.facebook.com/business/learn/certification}}
        {}
        {}

\cventry{2021年1月}
        {HubSpot Academy}
        {行銷自動化專家認證}
        {\url{https://academy.hubspot.com/}}
        {}
        {}
\section{出版刊物}

\cventry{2023年11月}
        {從數據到決策：零售品牌的全通路行銷實踐}
        {商業周刊}
        {\url{https://www.businessweekly.com.tw/}}
        {}
        {分享如何整合線上與線下數據，建立可衡量的全通路行銷策略，並以實際案例說明會員成長與營收提升的關鍵。}
\section{項目}

\cventry{2022年1月–2022年12月}
        {為台灣知名零售品牌重新定位，吸引 18-35 歲年輕客群。}
        {品牌年輕化再造計畫}
        {\url{https://www.yatinglin.com/projects/brand-revamp}}
        {}
        {\begin{itemize}
\item 進行消費者洞察與競品分析，重新定義品牌核心訊息與視覺風格
\item 規劃社群挑戰賽、KOL 合作與線下快閃店，創造 1,200 萬次社群曝光
\item 專案結束後，品牌在目標客群中的知名度提升 30\%
\end{itemize}
\textbf{關鍵字}：品牌定位、年輕化、社群行銷、消費者洞察}

\cventry{2023年3月–2023年11月}
        {運用 CRM 數據與自動化行銷，提升會員回購率與終身價值。}
        {會員忠誠度提升計畫}
        {\url{https://www.yatinglin.com/projects/loyalty-program}}
        {}
        {\begin{itemize}
\item 建立分群模型，針對不同消費週期設計個人化優惠與內容
\item 導入自動化 Email 與 LINE 訊息，開信率提升 45\%
\item 會員回購率在 9 個月內成長 28\%，帶動營收增加新台幣 3,200 萬元
\end{itemize}
\textbf{關鍵字}：CRM、自動化行銷、會員經營、數據分析}
\section{興趣愛好}

\cvline{藝文與設計}{當代藝術、平面設計、攝影}
\cvline{戶外活動}{登山、馬拉松、單車}
\cvline{閱讀與寫作}{商業書寫、行銷趨勢、專欄投稿}
\section{志願服務}

\cventry{2020年3月–2023年12月}
        {志工講師}
        {台灣行銷傳播專業認證協會}
        {}
        {}
        {\begin{itemize}
\item 擔任行銷認證課程志工講師，分享數位行銷與品牌策略實務
\item 協助協會舉辦產業論壇與學生工作坊，累計參與者超過 500 人
\item 參與偏鄉數位行銷教育計畫，協助在地小農建立社群媒體行銷能力
\end{itemize}}
    
\end{document}